\documentclass[12pt, letterpaper]{article}

% Essential Packages
\usepackage[utf8]{inputenc}
\usepackage[T1]{fontenc}
\usepackage{geometry}
\usepackage{setspace} % For line spacing
\usepackage{fancyhdr}
\usepackage{times} % Times New Roman
\usepackage[hidelinks]{hyperref} % For URLs, links hidden for print look
\usepackage[none]{hyphenat} % Prevent excessive hyphenation
\usepackage{titlesec} % For custom heading formatting

% Formatting Setup
\geometry{margin=1in}
\singlespacing % Set to Single Spacing per request

% Header Configuration
\pagestyle{fancy}
\fancyhf{}
\renewcommand{\headrulewidth}{0pt}
\fancyhead[R]{Balbale \thepage}

% Law Review Heading Configuration
% Level 1: Roman Numerals, Centered, Bold, ALL CAPS
\titleformat{\section}[block]{\bfseries\centering}{\Roman{section}.}{0.5em}{\MakeUppercase}

% Level 2: Letters, Italicized
\titleformat{\subsection}[block]{\itshape\centering}{\Alph{subsection}.}{0.5em}{}

\begin{document}

% Title Block
\begin{center}
    \textbf{\Large Embodied Harms and Inferred Data: Redefining Privacy in Virtual Reality}
    
    \vspace{1em}
    \textit{Sean Balbale}
\end{center}

\vspace{1em}

\section{Introduction}

The emergence of Extended Reality (XR),\footnote{Extended Reality (XR) is an umbrella term for immersive technologies that blend the physical and virtual worlds. It encompasses the entire spectrum, including Virtual Reality (VR), which fully immerses the user in a digital environment; Augmented Reality (AR), which overlays digital information onto the real world; and Mixed Reality (MR), which allows digital and physical objects to co-exist and interact in real-time.} an umbrella term for Virtual and Augmented Reality (VR/AR) technologies, represents a paradigm shift in human-computer interaction. Unlike previous digital eras where extensive data collection was technically optional, immersive technologies require the continuous collection of intimate user data to function. This new data-gathering model creates profound ethical and technical challenges. The primary risk is not identification—such as iris or fingerprint scans—but inference. These systems passively collect "behavioural biometrics," including subtle head motions, hand movements, and eye-gaze patterns.\footnote{Dilshani Kumarapeli et al., \textit{Privacy Threats of Behaviour Identity Detection in VR}, 5 \textsc{Frontiers in Virtual Reality} § 1 (2024).}

The power of these new data streams is significant; technical analyses demonstrate that VR users can be "uniquely recognised with approximately 90\% accuracy using head motions" alone.\footnote{\textit{Id.} § 2.1.} This data enables platforms to infer sensitive information that individuals did not choose to disclose, including mental and physical health status,\footnote{\textit{Id.} § 2.2.} sexual preferences,\footnote{Ellysse Dick, \textit{Balancing User Privacy and Innovation in Augmented and Virtual Reality}, \textsc{Info. Tech. \& Innovation Found.} 11 (Mar. 4, 2021) [hereinafter Dick].} and concentration levels.\footnote{Vasilis Xynogalas \& M.R. Leiser, \textit{The Metaverse: Searching for Compliance with the General Data Protection Regulation}, 14 \textsc{Int'l Data Privacy L.} 91 (2024).}

Consequently, the locus of data collection moves from an active process (such as clicks or text entries) to a passive one. This Article examines four primary dimensions of this shift: (1) competing architectures of control; (2) the nature of immersive harm; (3) the obsolescence of existing law; and (4) the search for a new governance model. A critical gap persists between voluntary, corporate-led technical solutions and the growing necessity for legally mandated public frameworks.


\end{document}